\sscp{Austronesian incursion/migration(s): \tsc{Timor-Babar}}{15-06-06}
It is not yet clear where Austronesian-speaking peoples came
from to start what eventually developed into the \tsc{Timor-Babar} subgroup (see \crf{ch:TB}).
Different patterns and discontinuities in the sound correspondences
do not favour a simple continuation from \tsc{Bima-Lembata}.
As mentioned in \srf{15-06-01},
it looks like they got off a different boat.
Whether they came directly or via an unknown intermediate node is not clear.
Nor does it appear there is a connection linking from central
Maluku as suggested by \citet{Blust1993,Blust2008}.
%See \srf{15-05-03}.

Within \tsc{Timor-Babar}, conservative features are found at both extremes.
In the far west (Helong),
and the far east (\tsc{Southwest Tanimbar}),
there are conservative retentions,
for example, in the labials.
There are also pockets of great diversity in many regions,
perhaps the most extreme being in the \tsc{Southwest Maluku}
node, which appears to show heavy substrate influence.
The direction of sound changes from east-to-west or west-to-east
within the subgroup is not clear,
for example with *p or *b.
There is also the appearance of some languages in the \tsc{Eastern
Timor} node having migrated either eastward or westward to their
current locations, and being in intense contact with the \tsc{Central Timor} subgroup
as well as the Papuan language of Bunak.

Given the complexities and evidence of prolonged
and intense contact throughout the entire subgroup,
the principle of “the area with the greatest diversity reflects
the homeland” should not be assumed.
Having said that, we note that the greatest diversity is found
in the eastern regions of the subgroup.